% cox and mullender, semaphores.
% 
% process has one thread with two stacks
%  
% pike et al, sleep and wakeup
\chapter{Scheduling}
\label{CH:SCHED}

Any operating system is likely to run with more processes than the
computer has CPUs, so a plan is needed to time-share the
CPUs among the processes. Ideally the sharing would be transparent to user
processes.  A common approach is to provide each process
with the illusion that it has its own virtual CPU by
\indextext{multiplexing}
the processes onto the hardware CPUs.
This chapter explains how xv6 achieves this multiplexing.
%% 
\section{Multiplexing}
%% 

Xv6 multiplexes by switching each CPU from one process to
another in two situations. First, xv6's
\lstinline{sleep}
and
\lstinline{wakeup}
mechanism switches when a process makes a system
call that blocks (has to wait), for example in \lstinline{read} or
\lstinline{wait}.
Second, xv6 periodically forces a switch to cope with
processes that compute for long periods without blocking.
The former are voluntary switches;
the latter are called involuntary.
This multiplexing creates the illusion that
each process has its own CPU.

Implementing multiplexing poses a few challenges. First, how to switch from one
process to another? 
The basic idea is to save and restore CPU registers,
though the fact that this cannot be expressed in C makes it tricky.
Second, how to force
switches in a way that is transparent to user processes?  Xv6 uses the
standard technique in which a hardware timer's interrupts drive context switches.
Third, all of the CPUs switch among the same set of processes, so a
locking plan is necessary to avoid mistakes such as two CPUs deciding
to run the same process at the same time. Fourth, a process's
memory and other resources must be freed when the process exits,
but it cannot do all of this itself
because (for example) it can't free its own kernel stack while still using it.
Fifth, each CPU of a multi-core machine must remember which
process it is executing so that system calls affect the correct
process's kernel state.
Finally,  \lstinline{sleep} and \lstinline{wakeup} allow a process to give up
the CPU and wait to be woken up by another process or interrupt.
Care is needed to avoid races that result in
the loss of wakeup notifications.

%% 
\section{Code: Context switching}
%% 

\begin{figure}[t]
\center
\includegraphics[scale=0.5]{fig/switch.pdf}
\caption{Switching from one user process to another.  In this example, xv6 runs with one CPU (and thus one scheduler thread).}
\label{fig:switch}
\end{figure}

Figure~\ref{fig:switch} 
outlines the steps involved in switching from one
user process to another:
a trap (system
call or interrupt) from the old process's user space to its kernel thread,
a context switch to the current CPU's scheduler thread, a context
switch to a new process's kernel thread, and a trap return
to the user-level process.
A CPU's scheduler thread is responsible for picking the next
process for that CPU to run, perhaps scanning repeatedly
(``idling'') if no process currently wants to run.
In this section we'll examine the mechanics of switching
between a kernel thread and a scheduler thread.

Switching from one thread to another involves saving the old thread's
CPU registers, and restoring the previously-saved registers of the
new thread; the fact that
the stack pointer and program counter
are saved and restored means that the CPU will switch stacks and
switch what code it is executing.

The function
\indexcode{swtch}
saves and restores registers for a kernel thread switch.
\lstinline{swtch}
doesn't directly know about threads; it just saves and
restores sets of RISC-V registers, called 
\indextext{contexts}.
When it is time for a process to give up the CPU,
the process's kernel thread calls
\lstinline{swtch}
to save its own context and restore the scheduler's context.
Each context is contained in a
\lstinline{struct}
\lstinline{context}
\lineref{kernel/proc.h:/^struct.context/},
itself contained in a process's
\lstinline{struct proc}
or a CPU's
\lstinline{struct cpu}.
\lstinline{swtch}
takes two arguments:
\indexcode{struct context}
\lstinline{*old}
and
\lstinline{struct}
\lstinline{context}
\lstinline{*new}.
It saves the current registers in
\lstinline{old},
loads registers from 
\lstinline{new},
and returns.

Let's follow a process through
\lstinline{swtch} 
into the scheduler.
We saw in Chapter~\ref{CH:TRAP}
that one possibility at the end of an interrupt is that 
\indexcode{usertrap}
calls 
\indexcode{yield}.
\lstinline{yield}
in turn calls
\indexcode{sched},
which calls
\indexcode{swtch}
to save the current context in
\lstinline{p->context}
and switch to the scheduler context previously saved in 
\indexcode{cpu->context}
\lineref{kernel/proc.c:/swtch..p/}.

\lstinline{swtch}
\lineref{kernel/swtch.S:/swtch/}
saves callee-saved registers but not caller-saved registers.
The RISC-V calling convention requires that if a function needs the
value in a caller-saved register, the compiler-generated instructions
for that function must save the register to the stack before calling
any other function, and restore from the stack when that other
function returns. So \lstinline{swtch} can rely on the function that
called it having already saved the caller-saved registers (either
that, or the calling function didn't need the values in the
registers).

\lstinline{swtch} knows the offset of each
register's field in 
\lstinline{struct context}.
It does not save the program counter.
Instead,
\lstinline{swtch}
saves the
\lstinline{ra} register,
which holds the return address from which
\lstinline{swtch}
was called.
Now
\lstinline{swtch}
restores registers from the new context,
which holds register values saved by a previous
\lstinline{swtch}.
When 
\lstinline{swtch}
returns, it returns to the instructions pointed to
by the restored
\lstinline{ra}
register, that is,
the instruction from which the new thread previously
called \lstinline{swtch}.
In addition, it returns on the new thread's stack,
since that's where the restored \lstinline{sp} points.

In our example, 
\indexcode{sched}
called
\indexcode{swtch}
to switch to
\indexcode{cpu->context},
the per-CPU scheduler context.
That context was saved at the point in the past when
\lstinline{scheduler}
called
\lstinline{swtch}
\lineref{kernel/proc.c:/swtch.&c/}
to switch to the process that's now giving up the CPU.
When the
\indexcode{swtch}
we have been tracing returns,
it returns not to
\lstinline{sched}
but to 
\indexcode{scheduler},
with the stack pointer in the current CPU's
scheduler stack.
%% 
\section{Code: Scheduling}
%% 

% why a separate thread

The last section looked at the low-level details of
\indexcode{swtch};
now let's take 
\lstinline{swtch}
as a given and examine 
switching from one process's kernel thread
through the scheduler to another process.
The scheduler exists in the form of a special thread per CPU, each running the
\lstinline{scheduler}
function.
This function is in charge of choosing which process to run next.
Each CPU has its own scheduler thread because more than one may be
looking for something to run at any given time. Process switching
always goes through the scheduler thread, rather than direct from one
process to another, to avoid some situations in which there would be
no stack on which to execute the scheduler (e.g. if the old process
has exited, or there is no other process that currently wants to run).

A process
that wants to give up the CPU must
acquire its own process lock
\indexcode{p->lock},
release any other locks it is holding,
update its own state
(\lstinline{p->state}),
and then call
\indexcode{sched}.
You can see this sequence in
\lstinline{yield}
\lineref{kernel/proc.c:/^yield/},
\texttt{sleep}
and
\texttt{exit}.
\lstinline{sched}
double-checks some of those requirements
\linerefs{kernel/proc.c:/if..holding/,/running/}
and then checks an implication:
since a lock is held, interrupts should be disabled.
Finally,
\indexcode{sched}
calls
\indexcode{swtch}
to save the current context in 
\lstinline{p->context}
and switch to the scheduler context in
\indexcode{cpu->context}.
\lstinline{swtch}
returns on the scheduler's stack
as though
\indexcode{scheduler}'s
\lstinline{swtch}
had returned
\lineref{kernel/proc.c:/swtch.*contex.*contex/}.
The scheduler continues its
\lstinline{for}
loop, finds a process to run, 
switches to it, and the cycle repeats.

We just saw that xv6 holds
\indexcode{p->lock}
across calls to
\lstinline{swtch}:
the caller of
\indexcode{swtch}
acquires the lock, but it's the switched-to thread
that releases the lock.
This arrangement is unusual: it's more common for
the thread that acquires a lock to also release it.
Xv6's context switching must break this convention because
\indexcode{p->lock}
protects invariants on the process's
\lstinline{state}
and
\lstinline{context}
fields that are not true while executing in
\lstinline{swtch}.
For example, if
\lstinline{p->lock}
were not held during
\indexcode{swtch},
a different CPU might decide
to run the process after 
\indexcode{yield}
had set its state to
\lstinline{RUNNABLE},
but before 
\lstinline{swtch}
caused it to stop using its own kernel stack.
The result would be two CPUs running on the same stack,
which would cause chaos.
Once \lstinline{yield} has started to modify a running process's state
to make it
\lstinline{RUNNABLE},
\lstinline{p->lock} must remain held until the invariants are restored:
the earliest correct release point is after
\lstinline{scheduler}
(running on its own stack)
clears
\lstinline{c->proc}.
Similarly, once 
\lstinline{scheduler}
starts to convert a \lstinline{RUNNABLE} process to
\lstinline{RUNNING},
the lock cannot be released until the process's kernel thread
is completely running (after the
\lstinline{swtch},
for example in
\lstinline{yield}).

The only place a kernel thread gives up its CPU is in
\lstinline{sched},
and it always switches to the same location in \lstinline{scheduler}, which
(almost) always switches to some kernel thread that previously called
\lstinline{sched}. 
Thus, if one were to print out the line numbers where xv6 switches
threads, one would observe the following simple pattern:
\lineref{kernel/proc.c:/swtch..c/},
\lineref{kernel/proc.c:/swtch..p/},
\lineref{kernel/proc.c:/swtch..c/},
\lineref{kernel/proc.c:/swtch..p/},
and so on.
Procedures that intentionally transfer control to each other via thread
switch are sometimes referred to as
\indextext{coroutines}; 
in this example,
\indexcode{sched}
and
\indexcode{scheduler}
are co-routines of each other.

There is one case when the scheduler's call to
\indexcode{swtch}
does not end up in
\indexcode{sched}.
\lstinline{allocproc} sets the context \lstinline{ra}
register of a new process to
\indexcode{forkret}
\lineref{kernel/proc.c:/^forkret/},
so that its first \lstinline{swtch} ``returns'' 
to the start of that function.
\lstinline{forkret}
exists to release the 
\indexcode{p->lock};
otherwise, since the new process needs to
return to user space as if returning from \lstinline{fork},
it could instead start at
\lstinline{usertrapret}.

\lstinline{scheduler}
\lineref{kernel/proc.c:/^scheduler/} 
runs a loop:
find a process to run, run it until it yields, repeat.
The scheduler
loops over the process table
looking for a runnable process, one that has
\lstinline{p->state} 
\lstinline{==}
\lstinline{RUNNABLE}.
Once it finds a process, it sets the per-CPU current process
variable
\lstinline{c->proc},
marks the process as
\lstinline{RUNNING},
and then calls
\indexcode{swtch}
to start running it
\linerefs{kernel/proc.c:/Switch.to/,/swtch/}.


%% 
\section{Code: mycpu and myproc}
%% 

Xv6 often needs a pointer to the current process's \lstinline{proc}
structure. On a uniprocessor one could have a global variable pointing
to the current \lstinline{proc}. This doesn't work on a multi-core
machine, since each CPU executes a different process. The way to
solve this problem is to exploit the fact that each CPU has its own
set of registers.

While a given CPU is executing in the kernel, xv6 ensures that the CPU's \lstinline{tp}
register always holds the CPU's hartid.
RISC-V numbers its CPUs, giving each
a unique \indextext{hartid}.
\indexcode{mycpu}
\lineref{kernel/proc.c:/^mycpu/}
uses \lstinline{tp} to index an array
of \lstinline{cpu} structures and return the one
for the current CPU.
A
\indexcode{struct cpu}
\lineref{kernel/proc.h:/^struct.cpu/}
holds a pointer to the \lstinline{proc}
structure of
the process currently running
on that CPU (if any),
saved registers for the CPU's scheduler thread,
and the count of nested spinlocks needed to manage
interrupt disabling.

Ensuring that a CPU's \lstinline{tp} holds the CPU's
hartid is a little involved, since user code is free
to modify \lstinline{tp}. \lstinline{start} sets the \lstinline{tp}
register early in the CPU's boot sequence, while still in machine mode
\lineref{kernel/start.c:/w_tp/}.
\lstinline{usertrapret} saves \lstinline{tp} in the trampoline
page, in case user code modifies it.
Finally, \lstinline{uservec} restores that saved \lstinline{tp}
when entering the kernel from user space
\lineref{kernel/trampoline.S:/make tp hold/}.
The compiler guarantees never to modify \lstinline{tp}
in kernel code.
It would be more convenient if xv6 could ask the RISC-V
hardware for the current hartid whenever needed,
but RISC-V allows that only in
machine mode, not in supervisor mode.

The return values of
\lstinline{cpuid}
and
\lstinline{mycpu}
are fragile: if the timer were to interrupt and cause
the thread to yield and later resume execution on a different CPU,
a previously returned value would no longer be correct.
To avoid this problem, xv6 requires that callers 
disable interrupts, and only enable
them after they finish using the returned
\lstinline{struct cpu}.

The function
\indexcode{myproc}
\lineref{kernel/proc.c:/^myproc/}
returns the
\lstinline{struct proc}
pointer
for the process that is running on the current CPU.
\lstinline{myproc}
disables interrupts, invokes
\lstinline{mycpu},
fetches the current process pointer
(\lstinline{c->proc})
out of the
\lstinline{struct cpu},
and then enables interrupts.
The return value of
\lstinline{myproc}
is safe to use even if interrupts are enabled:
if a timer interrupt moves the calling process to a
different CPU, its
\lstinline{struct proc}
pointer will stay the same.


%% 
\section{Real world}
%% 

The xv6 scheduler implements a simple scheduling policy, which runs each process
in turn.  This policy is called
\indextext{round robin}.
Real operating systems implement more sophisticated policies that, for example,
allow processes to have priorities.  The idea is that a runnable high-priority process
will be preferred by the scheduler over a runnable low-priority process.   These
policies can become complex quickly because there are often competing goals: for
example, the operating system might also want to guarantee fairness and
high throughput.

%% 
\section{Exercises}
%% 

\begin{enumerate}

\item Modify xv6 to use only one context switch when switching from
  one process's kernel thread to another, rather than switching
  through the scheduler thread. The yielding thread will need to select
  the next thread itself and call \texttt{swtch}. The challenges will
  be to prevent multiple CPUs from executing the same thread accidentally;
  to get the locking right; and to avoid deadlocks.

\end{enumerate}
